% MT036 / explanation — canonical editable content.
这里输入是68维音频特征，输出是经纬度。10,597条记录分布不均，还含异常值；要比较的是模型怎样表示规律、计算成本和数据缺陷会怎样影响预测。


\textbf{GP的优势是带模型不确定性的预测。} 标准GP给出预测均值及方差，除了位置估计，还可以提醒某些输入缺乏相近训练证据。不同地区样本不均时，这种信息有用；但可信程度仍依赖核与噪声模型，不是样本少就必然得到准确的不确定性。


\textbf{GP的代价是精确稠密计算较重。} \(N=10{,}597\)的协方差矩阵有112,296,409个元素，若每个8字节，单矩阵约0.90 GB（十进制），还没有计算分解和其他中间量。标准精确求解也有立方级时间增长。高斯观测模型对极端目标偏差可能敏感，需要核查异常值或采用合适的稳健扩展。


\textbf{RF能表达非线性及特征交互，通常较容易扩展训练。} 不同树可独立拟合、并行预测。但普通回归树叶子常输出训练目标均值，极端经纬度标签仍可能把均值拉偏；不能称其对异常值免疫，也不能凭名称断言它必然比GP更敏感。


两个方法都不会自动补齐少样本地区。验证应按实际泛化目标设计，并分别查看各地区误差，防止多数地区的平均成绩掩盖少数地区失败。普通RF也不会自动提供校准过的预测概率分布；可以构造不确定性估计，但应另行验证。


\textbf{出处：} 2021B*期中 Q10，原卷第4页（10分）


\textbf{答案依据：} 2021B*期中 Q10 随卷／配套参考答案；条件与解释经复核。
